\documentclass[../veristore_security_model.tex]{subfiles}

\begin{document}
\section{Dependencies}

% COMMENT: Enumerate all external libraries, systems, and services that veri-store relies upon.
% For each dependency, note what functionality it provides and what security assumptions you're
% making about it. Examples: Python standard library, galois/pyfinite for finite field arithmetic,
% numpy for array operations, network stack (TCP/IP), file system for local storage.

\subsection{External Libraries}

\subsubsection*{Mathematical operations and erasure coding}
\begin{itemize}
    \item \texttt{\textbf{galois} [$\geq$ v0.3.8]}: Finite field arithmetic over $\mathbb{F}_{256}$, primarily utilized in developing Cauchy matrices
    \item \texttt{\textbf{numpy} [$\geq$ v1.26.0]}: Array operations and polynomial manipulation
    \item \texttt{\textbf{reedsolo} [$\geq$ v1.7.0]} (\textit{optional}): Encoding and decoding data before dispersal
\end{itemize}

\subsubsection*{Web framework}
\begin{itemize}
    \item \texttt{\textbf{fastapi} [$\geq$ v0.110.0]}: REST API for client-server communication
    \item \texttt{\textbf{uvicorn} [$\geq$ v0.29.0]}: ASGI server for hosting the API
\end{itemize}

\subsubsection*{HTTP client}
\begin{itemize}
    \item \texttt{\textbf{httpx} [$\geq$ v0.27.0]}: For inter-server communication and fragment distribution
\end{itemize}

\subsubsection*{API data validation}
\begin{itemize}
    \item \texttt{\textbf{pydantic} [$\geq$ v2.6.0]}: Data validation and parsing for API requests/responses
\end{itemize}

\subsubsection*{Testing}
\begin{itemize}
    \item \texttt{\textbf{pytest} [$\geq$ v8.0.0]}: For running unit and integration tests
    \begin{itemize}
        \item \texttt{\textbf{pytest-asyncio} [$\geq$ v0.23.0]}: For testing asynchronous code paths
    \end{itemize}
\end{itemize}

\subsection{System Dependencies}
Veri-store also depends on the underlying platform. We treat these as given and do not try to secure them ourselves:
\begin{itemize}
    \item \textbf{Network stack}: We assume a working TCP/IP stack that can deliver HTTP requests between clients and servers. Attacks on routing or congestion control are out of scope.
    \item \textbf{File system}: We rely on the local file system to store fragment files and metadata. We assume basic guarantees such as ``writes that complete are durable'' and that file permissions are configured correctly by the operator.
    \item \textbf{Operating system}: We assume the OS provides process isolation, memory protection, and a reasonably bug-free implementation of standard syscalls. Kernel-level exploits and privilege escalation are not modeled here.
\end{itemize}

\newpage
\end{document}